\section{Part P: the opened literature and the comparison maps}

\begin{result}{P.1}{Known}{Literature audit, with the scope of the maps stated}
The opened sources give the following results. The claims audited in Part A must be compared with these exact statements.
\end{result}

\paragraph{Genus zero and the two completions.}
For a topological surface $S_{g,n}$, let $\Gamma_{g,n}$ be its pure orientation-preserving mapping class group. We write
\[
 \widehat\Gamma_{g,n}=\varprojlim_{K\triangleleft_f\Gamma_{g,n}}
          \Gamma_{g,n}/K,
 \qquad
 \check\Gamma_{g,n}
 =\overline{\Gamma_{g,n}}^{\,\Out(\widehat\pi_1(S_{g,n}))}.
\]
The continuous homomorphism
$\widehat\Gamma_{g,n}\twoheadrightarrow\check\Gamma_{g,n}$
is induced by the discrete outer action on the surface fundamental group. It is surjective because its image is compact and contains the dense discrete image. Its kernel is exactly the congruence kernel for that action. Every passage between these two groups below uses this homomorphism.

Harbater--Schneps \cite[Main Theorem and its corollary]{HS2000} identifies the subgroup of outer automorphisms satisfying their inertia and permutation conditions with $\GTz$ at four points and with $\GT$ at five or more points. Their construction uses the actual forgetting homomorphisms. Hoshi--Minamide--Mochizuki \cite[Corollary 2.8]{HMM2022} removes the inertia restrictions: for $n\geq5$,
\[
 \Out(\widehat\Gamma_{0,n})\simeq\mathfrak S_n\times\GT.
 \tag{P.1}
\]
To translate their notation, take the configuration space of $n-3$ ordered points on $\mathbb P^1\setminus\{0,1,\infty\}$; its fundamental group is $\Gamma_{0,n}$. Their $\Pi_{n-3}$ with $\Sigma$ the set of all primes is its full profinite completion. The two factors in (P.1) are respectively relabelling and the automorphisms commuting with all the named forgetting maps. This gives the precise input used at eight points in Part A.

Boggi \cite[Corollary 5.4 and Theorem 1.4(ii)]{CA_Boggi2020} supplies congruence faithfulness in genus zero and in the hyperelliptic cases containing genus one and two. In genus zero the reduction is induction through the Birman sequences, starting with the thrice-punctured sphere. For genus one, his notation $\Upsilon_{1,n}=\Gamma_{1,n}$ and the level-two subgroup in Theorem 1.4 have finite index; its faithful completed action therefore gives the full-group comparison, since the remaining finite quotient is also detected. We use this result only in these stated cases. In higher genus the displayed epimorphism remains the comparison map; we do not insert an unproved inverse.

\paragraph{The genus-one relation.}
Hatcher--Lochak--Schneps \cite{ET_HLS2000} proves the presentation of the Teichm\"uller tower by low-complexity moves and relations. Nakamura--Schneps \cite[Definition 1.1, Theorems 1.2--1.4, and Sections 7--9]{CA_NS2000} constructs a tower action for the subgroup defined there by the additional genus-one relations. In their proof, the mixed six-move relation is an actual relation on the twice-punctured torus. The extension to all of $\GT$ asserted by EXIT\_TOWER requires a derivation of that relation from the Drinfeld relations; quoting this subgroup action does not itself supply that derivation.

Lochak \cite[Definition 5.4 through Proposition 5.9]{Lochak2025Aut} considers pairs of parameters, with their mixed compatibility equation. The first projection is injective and its image is described by the existence of a second GT parameter. EXIT\_TOWER (E.200) reproduces this equation, with the same placements of the powers and of the two parameters. Its further assertions are that the projection is onto and that both parameters agree. Those assertions require the ten-point comparison checked below. Lochak's procongruence curve-complex conclusions use $\check\Gamma$; they do not provide the missing inverse to $\widehat\Gamma\twoheadrightarrow\check\Gamma$ in higher genus.

The later opened papers of Lochak--Nakamura--Schneps \cite{LNS2026} and Basualdo Bonatto--Robertson \cite{ET_BonattoRobertson2026} also state their own precise scopes. The former constructs the indicated actions on the groups with zero or one marked point. The latter's Theorem 4.9 uses the Nakamura--Schneps subgroup, while its Proposition 4.7 identifies the genus-zero automorphism group with $\GT$. These results do not replace the mixed-relation verification in our proof.

\paragraph{Profinite and prounipotent comparisons.}
Drinfeld \cite[Section 4, printed pp.~844--847]{Drinfeld1991} defines the profinite, pro-$\ell$, and characteristic-zero prounipotent versions and gives the homomorphisms
\[
 \GT\longrightarrow\mathrm{GT}_{\ell}
       \longrightarrow\mathrm{GT}(\mathbb Q_\ell).
 \tag{P.2}
\]
The homomorphisms come from the free-group completion maps and retain the two-cycle, three-cycle, and pentagon equations by substitution. To transport a counterexample in the last group back through (P.2), one must exhibit a preimage satisfying the full profinite gentle conditions. Merely having the homomorphisms proves no such lifting assertion. We inspected the defining pages directly; they state the role of the pentagon in genus-zero extension, and do not establish a strict profinite inclusion by that argument. The supplied refereed LM7\_PUSH chain gives a continuous quadratic character on $\GTz$ and its arithmetic vanishing. Its vanishing on $\GT$ is the separate question of Part A.

\paragraph{Finite subgroups.}
Boggi--Zalesskii \cite[Theorem A, equivalently Theorem 4.3]{CA_BZ2025} gives conjugacy and normalizer/centralizer density for a single finite $p$-subgroup under its cohomological hypotheses. It directly gives
\[
 N_{\widehat M}(H)=\overline{N_M(H)},\qquad
 C_{\widehat M}(H)=\overline{C_M(H)}
 \tag{P.3}
\]
for the finite $p$-subgroups of the low-genus mixed groups used below. The derivation of the hypotheses and of the full topology of the particular discrete centralizers is supplied in Part A; density alone is not substituted for that topology assertion.

At the end of the proof of Boggi--Zalesskii Theorem 5.1(iii), and in Boggi \cite[Proposition 2.4]{CA_Boggi2026}, the equality between an intersection of closures and the closure of the simultaneous centralizer needs a separate argument. The displayed proof does not provide it for arbitrary finite subgroups. This observation does not constitute a counterexample to either theorem. The finite $p$-group applications in this paper use (P.3) directly. The solvable case is recovered from normalizer density by the finite conjugation map to $\Aut(H)$, as proved below.

\paragraph{Search scope and attribution.}
We searched the named authors' primary publication pages, arXiv records and sources, and the primary papers cited above, including the 2026 papers just described. We found no earlier proof in those opened sources of the complete chain asserted by EXIT\_TOWER E.39--E.44 and E.49. This records the scope of the search. It is not a claim of priority or a theorem about the entire literature. The supplied returns are evidence to be audited, and are cited throughout by their run-prefixed identifiers.
